\section{Post-Training}
\label{sec:post-training}

\subsection{Method}

Our post-training pipeline follows a three-stage paradigm:
initializing baseline agent capabilities via supervised fine-tuning (SFT),
developing specialized domain experts at varying reasoning effort via Reinforcement Learning (RL),
and consolidating these domain-specific policies into a single model using Multi-Teacher On-Policy Distillation (MOPD).

\subsubsection{Supervised Fine-Tuning}
\label{sec:post-sft}
The SFT stage establishes a high-quality cold-start policy for the subsequent RL stage.
Building on the SFT pipeline of previous Kimi models \cite{kimiteam2025kimik2openagentic,kimik25}, we expand the SFT dataset for \kimi{3}, substantially broadening its coverage of complex agentic tasks.
Specifically, we synthesize data trajectories using domain-specialized models from the prior Kimi series, followed by multi-stage verification and human-in-the-loop annotation.
To represent these complex agentic trajectories consistently, we serialize all data with our XTML-based chat template (eXtensible Token Markup Language; see \S~\ref{sec:post-chat-template} for details).
Collectively, these steps yield a large-scale instruction dataset that endows \kimi{3} with adaptive reasoning, precise tool calling, and robust execution in long-horizon agentic scenarios.
In addition, we apply quantization-aware training (QAT) from the SFT stage onward, with MXFP4 weights and MXFP8 activations (\S~\ref{sec:mxfp4-qat}).

\subsubsection{Reinforcement Learning}
\label{sec:post-rl}

\begin{figure}[t]
    \centering
    \includegraphics[width=0.8\textwidth]{figures/step_scaling_merged.pdf}
    \caption{Scores and the average assistant steps across a variety of public and in-house evaluations during RL. By scaling RL FLOPs, tool-call steps scale up consistently, accompanied by a comprehensive improvement in the model's overall capability.}
    \label{fig:step_scaling_merged_curve}
\end{figure}


While SFT provides a solid cold-start foundation, RL is critical to unlocking higher-order reasoning and execution capabilities. 
Rather than training specialized RL models for individual tasks, we scale RL across three broad domains, each encompassing a wide spectrum of sub-tasks, and train a single expert for each domain at every reasoning effort level:
(i) \emph{general tasks}, spanning general experience, vision, reasoning, faithfulness, search capabilities, and knowledge work tasks; 
(ii) \emph{general agents}, spanning long-horizon assistant tasks, deep research, and paragraph-level writing; and 
(iii) \emph{coding agents}, spanning software engineering (SWE), coding experience, kernel tasks, and web development. 
As shown in Figure~\ref{fig:step_scaling_merged_curve}, 
scaling RL FLOPs consistently improves a variety of capabilities across knowledge, reasoning, vision, general agent, and coding. 
Crossing these three domain experts with three reasoning effort levels in $\{\text{low}, \text{high}, \text{max}\}$ yields a total of nine expert models.


\paragraph{Algorithm}

To mitigate the long-tail latency that intensifies in long-horizon tasks, we extend the \emph{partial rollout} scheme from our synchronous RL framework \citep{kimik15,kimik25}. 
During the rollout phase of each iteration, we sample $K$ completions for each of $N$ prompts, maintaining an active workload of $N \times K$ trajectories. Rather than waiting for all rollouts to terminate, the generation phase pauses as soon as a fraction $\lambda \in (0, 1)$ of trajectories completes (i.e., $\lambda NK$), allowing policy optimization to proceed without execution stragglers. Paused rollouts are enqueued and prioritized for resumption at the start of the next iteration, powered by our sandbox infrastructure (\S~\ref{sec:sandbox}). 
Once all $K$ responses for a prompt complete, they are immediately dispatched for policy optimization, which follows the algorithm in \kimi{2.5} \citep{kimik25}. 
Under our partial rollout scheme, an individual long-horizon trajectory naturally spans multiple iterations, introducing data staleness that threatens training stability. 
Our policy optimization algorithm inherently tolerates such an extreme off-policy regime through a per-token regularization.  By constraining policy updates within a localized neighborhood, this regularization enables the algorithm to robustly handle highly stale data and sustains training stability.



\paragraph{Reasoning Effort RL}

To fine-tune reasoning effort while maximizing token efficiency, we implement a per-problem budget control mechanism during RL \citep{kimik25}. 
We associate each problem $x$ with an initial token budget $b_0(x)$ estimated from the cold-start model, and override the task reward with $-1$ for trajectories whose total token budget $T(y)$ exceeds a scaled threshold $\tau \cdot b_0(x)$. 
For general tasks, $T(y)$ measures the number of thinking tokens, whereas for agentic tasks, $T(y)$ accounts for the cumulative output tokens, including both reasoning traces and tool-call arguments. 
Training follows a stage-wise curriculum over the budget multiplier $\tau$. We first train a \emph{max-budget} variant with a relatively large $\tau$, while still capping the maximum budget to suppress excessive overthinking. We then anneal $\tau$ to smaller values to obtain the \emph{high}- and \emph{low}-effort expert models. 
The adjustment of $\tau$ is configured per domain under human-in-the-loop guidance. Trajectories produced by the resulting experts at all reasoning levels are jointly collected for supervised fine-tuning and multi-teacher on-policy distillation.

\iffalse
\begin{align}
r(x, y) = \begin{cases} -1, & T(y) > \tau \cdot b_0(x) \\[2pt] R_{\mathrm{task}}(y \mid x), & \text{otherwise} \end{cases} \, ,
\end{align}
\fi

\paragraph{Agentic Generative Reward Model}
For non-verifiable general tasks, we adopt an Agentic Generative Reward Model (GRM), retaining the tournament-style group reward with binary comparisons as in \kimi{2.5}~\cite{kimiteam2025kimik2openagentic,kimik25}. Beyond generic agentic capabilities for enhanced judgment, the agentic judge is required to follow a mandatory protocol: (1) read the outcome, product, or text output; (2) generate a rubric; (3) score each candidate against the rubric; and (4) record the rubric-assigned scores in a scorepad. To mitigate reward hacking toward increasingly verbose outputs, we apply a budget-based verbosity control analogous to the reasoning-effort control above: given an initial verbosity $\ell_0$ estimated from the cold-start model and a multiplier $\sigma$, a candidate whose output length exceeds $\sigma \cdot \ell_0$ automatically loses the binary comparison.





\subsubsection{Multi-Teacher On-Policy Distillation}

We adopt Multi-Teacher On-Policy Distillation (MOPD) to consolidate these domain-specialized capabilities across varying reasoning efforts into a unified model \citep{lu2025onpolicydistillation, xiao2026mimov2flash, deepseekv4}. 
During training, for a given domain $d$ and a sampled reasoning effort level $e \in \{\text{low}, \text{high}, \text{max}\}$, optimization is guided by the corresponding teacher model $\pi_{\text{teacher}}^{(d,e)}$ among the nine experts. 
Given an input query $x$ and the prefix response $y_{<t}$, the per-token OPD reward evaluated on $y_t$ between the teacher $\pi_{\text{teacher}}^{(d,e)}$ and the student $\pi_{\theta}$ is defined as:
\begin{align}
r^{d}_{\mathrm{opd}}(y_t \mid e, x, y_{<t}) = \mathrm{clip}\left( \mathrm{sg} \left(\log \frac{\pi_{\text{teacher}}^{(d,e)}(y_t \mid x, y_{<t})}{\pi_{\theta}(y_t \mid e, x, y_{<t})} \right) , -R_{\max}, R_{\max} \right)\, ,
\end{align}
where $\operatorname{sg}(\cdot)$ denotes the stop-gradient operator, and $R_{\max} > 0$ is a clipping threshold to constrain extreme advantage signals, thereby stabilizing RL training. 
This dense reward signal seamlessly integrates into our RL framework, naturally enabling infrastructure-level optimizations such as partial rollout training for long-horizon tasks. 
While we also experimented with more fine-grained top-$k$ distillation objectives, we observed no clear advantage in either convergence speed or final performance in our setting. 

\subsubsection{Deployment-Aware Post-Training}


\paragraph{MXFP4 Quantization-Aware Post-Training}
\label{sec:mxfp4-qat}

To reduce memory footprint and serving cost at deployment, we quantize the MoE expert weights --- which dominate the model's parameter memory --- to MXFP4~\citep{rouhani2023microscaling}, with activations computed in MXFP8, while all non-expert components (attention projections, latent MoE projections, shared experts, and MoE routers) remain in higher precision.
We perform quantization-aware training (QAT)~\citep{jacob2018quantization} throughout the entire post-training stage, covering both SFT and RL, so that the model adapts to quantization-induced precision loss.
During RL, rollout and training share the same quantization scheme --- eliminating the train--inference mismatch.


\paragraph{Draft Model Fine-Tuning}
\label{sec:post-eagle3}
Optimizing inference efficiency is crucial for serving complex, long-horizon agentic models.
\kimi{3} is pre-trained with a multi-token-prediction (MTP) layer that mirrors the structure of a backbone block.
As the draft model of EAGLE-3~\citep{li2025eagle3} comprises a single decoder layer whose structure matches the MTP layer, we fine-tune the pre-trained MTP layer into an EAGLE-3-style draft model, with the target model frozen and only the draft layer and its feature-fusion projection updated.
Following the training-time test protocol of EAGLE-3, the draft is unrolled for seven steps during training; beyond the first step, where the target-side features of the newest position are unavailable, the draft consumes its own outputs from earlier steps, mirroring the recurrent drafting procedure at inference.

The draft input fuses low-, mid-, and high-level features of the target model, taken from the outputs of the 1st, 4th, and final AttnRes blocks, respectively (\S~\ref{sec:attnres}).
These features are concatenated and projected to the hidden size by a bias-free matrix $\bm{W}_{\mathrm{E3}}$, initialized as $[\,\bm{0}\;\;\bm{0}\;\;\bm{I}\,]$ so that the fused representation coincides at initialization with the high-level feature $\bm{h}_{h}$ --- the input on which the MTP layer was pre-trained --- and gradually learns to incorporate the low- and mid-level features during fine-tuning.

The speedup of speculative decoding is governed by the per-token acceptance rate $\sum_{x \in \mathcal{V}} \min\!\left(p(x), q(x)\right)$ under lossless speculative sampling, where $p$ and $q$ denote the next-token distributions of the target and draft models.
Since minimizing the conventional KL-divergence surrogate does not guarantee maximizing this rate for a capacity-limited draft model, we directly optimize the likelihood-based LK loss~\citep{samarin2026lklosses}, the negative logarithm of the acceptance rate itself,
\begin{equation}
    \mathcal{L}_{\mathrm{LK}}
    = -\log \sum_{x \in \mathcal{V}}
    \min\!\left(p(x), q(x)\right),
\end{equation}
with $p$ and $q$ evaluated at temperature~1 and no auxiliary ground-truth cross-entropy term.
Draft fine-tuning follows the post-training QAT configuration (\S~\ref{sec:mxfp4-qat}), with MoE expert weights in MXFP4 and their input activations in MXFP8, while non-expert modules remain in higher precision.



\subsection{RL Task Synthesis and Agentic Environments}
\label{sec:post-envs}

The effectiveness of our RL framework relies heavily on rich, diverse, and robustly verifiable environments. To support scalable training across complex long-horizon tasks, we design a series of specialized white-box environments and task synthesis paradigms.

\subsubsection{Unified White-Box RL Environment}

Training with a single fixed agent harness can cause a model to overfit to a particular tool schema, system prompt, context management mechanism, or interaction protocol. To address this, we develop a unified white-box RL environment that represents an agent harness as a collection of configurable, composable modules, including tool interfaces, system prompts, context management strategies, skills, memories, subagents, and other components. Composing these modules through configuration, the environment can instantiate mainstream harnesses such as Kimi Code~\citep{kimicode}, Claude Code~\citep{claudecode}, Codex~\citep{openaicodex}, OpenClaw~\citep{openclaw}, and Hermes~\citep{hermesagent}, as well as entirely new ones. During RL training, we dynamically construct different harness configurations for different task groups,  exposing \kimi{3} to diverse combinations of these modules rather than the conventions of any single harness. The same abstraction also readily supports RL across various task domains, providing a scalable foundation for training more general-purpose agents.

\subsubsection{Knowledge-Graph-Guided Task Synthesis}
\label{sec:post-data}

\paragraph{Motivation and overview}
The quality and diversity of post-training tasks are largely determined by their source materials. Retrieval guided by fine-grained concepts surfaces specialized and underrepresented knowledge, while sampling across diverse concepts broadens domain coverage. To control both granularity and coverage at scale, we build a self-evolving, hierarchically organized knowledge graph that agents continuously expand through web-scale exploration across knowledge-intensive and coding domains. Figure~\ref{fig:post-data-pipeline} illustrates the task synthesis pipeline.

\begin{figure}[htb]
    \centering
    \includegraphics[width=0.9\textwidth]{figures/k3_task_pipeline_v76.pdf}
    \caption{Overview of knowledge-graph-guided task synthesis. The hierarchically organized knowledge graph represents concepts at multiple levels, ranging from broad domains to fine-grained concepts. Related nodes are sampled to form a keyword set that guides the retrieval of publicly available source materials. For each synthesis instance, the system selects a task type and uses the retrieved materials to synthesize a corresponding task.}
    \label{fig:post-data-pipeline}
\end{figure}

\paragraph{Agentic knowledge graph construction}
We construct the knowledge graph as a directed acyclic graph through recursive, agent-driven expansion. The expansion process begins with a predefined set of coarse-grained seed nodes. An agent instance is then assigned to each node and performs multiple web searches to investigate the corresponding concept. Before adding new nodes, the agent explores the existing graph to identify equivalent or related concepts, reuse existing nodes where appropriate, and minimize duplication. Edges are always directed from the coarser concept to the finer one, regardless of which endpoint the agent discovers first. Newly added nodes are subsequently assigned to agents for further exploration. A branch stops expanding when the assigned agent determines that the current concept is sufficiently atomic.

\paragraph{Material retrieval and task synthesis}
To target a desired distribution across domains and task types, the system samples nodes at varying levels of granularity, either individually or in related combinations. Keywords derived from the sampled nodes are combined with contextual information from their ancestors in the knowledge graph to formulate web queries. The retrieved real-world materials are assembled so that a synthesis agent produces training tasks of various task types.

\subsubsection{Verifiable Problems in Agentic Environments}
\label{sec:agentic-verifiable-problems}

We train \kimi{3} on verifiable problems in agentic environments; representative examples include multi-step complex information searching, where the model plans its research, gathers evidence from the web step by step, and produces a verifiable answer; the real day-to-day work of professionals, such as investment banking, data analysis, and legal practice, where the model decomposes a complex request, operates domain tools in a sandbox, and completes a deliverable over dozens to hundreds of steps; and multi-step verifiable visual reasoning over STEM problems, visual puzzles, and chart understanding. Each visual-reasoning trajectory is generated in an agent environment equipped with a Python interpreter in an isolated sandbox: the model iteratively writes and executes code to crop, zoom, or transform the input image, perform precise computation, or verify intermediate results, and receives the execution outputs --- including generated images --- as new observations over multiple interaction steps. As the model learns to perform more image operations and collect more observations, its performance on complex visual reasoning tasks steadily improves.


\subsubsection{Kernel Optimization Tasks}
\label{sec:post-envs-kernel}

To strengthen \kimi{3}'s GPU kernel optimization capabilities, we build a large-scale suite of kernel tasks ranging from single-operator kernels to fused mega-kernels, sourced from high-quality GitHub repositories such as Flash Linear Attention~\cite{yang2024fla}. The suite spans diverse GPU programming approaches, such as CUDA, Triton, CuTe DSL, Gluon, ThunderKittens~\cite{spector2025thunderkittens}, and TileLang~\cite{wang2025tilelang}, and covers widely used GPU architectures and numerical formats including BF16, FP8, and FP4. Rewards evaluate both correctness and performance: each kernel provides a PyTorch reference implementation, and solutions exceeding a predefined numerical error threshold receive zero reward. Performance is scored against an expert implementation, where matching it yields a reward of 0.5 and approaching the hardware roofline increases the reward toward 1. To ensure that rewards reflect genuine optimization, we develop a hacking-detection system that penalizes reward-hacking strategies such as CUDA graph replay, input caching, and precision reduction, and we continuously extend it with new safeguards as new hacking strategies are observed during \kimi{3}'s development.

\subsubsection{Personal Assistant Tasks}
\label{sec:week-task}
For long-horizon personal assistant tasks, we develop realistic mock implementations of widely used applications, such as Gmail, Notion, Slack, and Canvas. They preserve the core semantics of their real-world counterparts while enabling reproducible, large-scale interaction without external APIs or rate limits. Building on these mock applications, we design complex tasks inspired by real-world professional workflows in scenarios like human resources, legal services, and finance. In each task, the agent operates in a persistent, evolving environment over multiple simulated days and encounters dozens of interdependent events distributed across applications. A single rollout may involve up to thousands of tool calls and millions of context tokens. Each event carries its own evaluation criterion, assessed by deterministic rules or LLM-based evaluators. The initial workspace is constructed by agents that autonomously search the web for reference materials and transform them into a coherent, task-relevant environment. We also extend our RL framework to support such living environments, modeling complex event streams and the induced world-state transitions.



\subsubsection{Autonomous Execution Tasks}
\label{sec:post-envs-autonomous}

We introduce Autonomous Execution Tasks (AET), an environment paradigm that trains long-horizon agent intelligence through verify-in-the-loop optimization. Each task specifies an initial state, a constrained goal, a tool-based action space, execution budgets, and an independent verifier. Agents see only the objective, context, constraints, and verification interfaces, without reference trajectories or predefined procedures, and must autonomously perform task decomposition, tool selection, planning, error recovery, and termination. Rewards are grounded in the verifier's evaluation of the final environment state rather than the agent's self-reported completion. We design multiple types of verifiers that support diverse environments, including black-box system replication (Figure~\ref{fig:aet_blackbox_curve}), quantitative factor discovery, and tax auditing. In each environment, agents iteratively submit solutions, receive verifier feedback, and refine their strategies, training a general loop of hypothesizing, acting, analyzing feedback, and adapting. Reward hacking is mitigated by isolating agents from verifiers, pairing public verifiers that offer diagnostic feedback with hidden verifiers that evaluate held-out scenarios, and applying penalty-based rewards under limited submission budgets.











\begin{figure}[t]
    \centering
    \includegraphics[width=0.65\textwidth]{figures/AET_blackbox_curve.pdf}
    \caption{Completion curves on Camera Repair Management System, a black-box system replication task in which the agent reconstructs a hidden 3D-camera repair system as a web application through oracle queries. Completion denotes verifier-assessed task progress.}
    \label{fig:aet_blackbox_curve}
\end{figure}

\subsubsection{Web Development Tasks}
\label{sec:webdev-tasks-autonomous}

We construct a diverse suite of expert-curated web development tasks covering typical scenarios. Inputs range from one-line scene descriptions to multi-paragraph specifications; artifacts span websites, interactive games, 3D/WebGL scenes, data visualization, SVGs, and full-stack applications. Every task runs in a containerized sandbox and is rolled out under diverse agent scaffolds rather than a single fixed harness, to promote cross-scaffold generalization. Rewards consist of two components: deterministic checks and model judging by an internal reward model. Deterministic checks functionally test application behavior, and score structural and pixel-level similarity for tasks that replicate a reference. The reward is zeroed when a project fails to build, runs with errors, or fakes rather than implements the artifact. Model judging uses other models to perform source code inspection or to look at and interact with the output artifact. 




















